\uuid{6qrM}
\titre{Système linéaire/ Système échelonné}
\auteur{hanani ; mzari}
\organisation{exo7}
\licence{CC BY-NC-SA 4.0}
\datecreate{2020-01-20}
\niveau{L1}
\module{Algèbre}
\chapitre{Déterminant, système linéaire}
\sousChapitre{Système linéaire, rang}
\difficulte{3}
\theme{Systèmes d'équations linéaires, Système linéaire, Système échelonné}
\exo7qcmid{lille-2.systemes.021}

\contenu{
\begin{question}
\qclassification{200.03}
\qtags{motcle=Système linéaire/ Système échelonné}
Soit $(\mathtt{S})$ un système d'équations linéaires et $(\mathtt{S}_E)$ un système échelonné obtenu par la méthode de résolution du pivot de Gauss. Quelles sont les assertions vraies ?
\begin{answers}
\bad{$(\mathtt{S})$ admet une infinité de solutions si et seulement si toute équation de $(\mathtt{S}_E)$ dont le premier membre est nul a aussi son second membre nul.}
\good{$(\mathtt{S})$ n'admet pas de solution si et seulement s'il existe une équation de $(\mathtt{S}_E)$ ayant un premier membre nul et un second membre non nul.}
\bad{$(\mathtt{S})$ admet une unique solution si et seulement si le nombre d'équations de $(\mathtt{S}_E)$ dont le premier membre est non nul est égal au nombre d'inconnues.}
\good{Si le nombre d'équations de $(\mathtt{S}_E)$ dont le premier membre est non nul est strictement inférieur au nombre d'inconnues et les équations ayant un premier membre nul admettent aussi le second membre nul, alors $(\mathtt{S})$ admet une infinité de solutions.}
\end{answers}
\begin{explanations} D'après la méthode du pivot de Gauss, $(\mathtt{S})$ admet une solution si et seulement si toute équation de $(\mathtt{S}_E)$ dont le premier membre est nul a aussi son second membre nul.
\end{explanations}
\end{question}
}
